\documentclass[12pt,a4paper]{article}
\usepackage[utf8]{inputenc}
\usepackage[T1]{fontenc}
\usepackage[spanish]{babel}
\usepackage{geometry}
\usepackage{booktabs}
\usepackage{longtable}
\usepackage{array}
\usepackage{xcolor}
\usepackage{hyperref}
\usepackage{fancyhdr}
\usepackage{graphicx}
\usepackage{titlesec}
\usepackage{enumitem}
\usepackage{mdframed}
\usepackage{tabularx}
\usepackage{multirow}
\usepackage{amsmath}
\usepackage{amssymb}

\geometry{margin=2.5cm}

% ---- Colores del sistema ----
\definecolor{primaryblue}{RGB}{33,150,243}
\definecolor{criticalred}{RGB}{244,67,54}
\definecolor{normalgreen}{RGB}{76,175,80}
\definecolor{warningorange}{RGB}{255,152,0}
\definecolor{codegray}{RGB}{245,245,245}
\definecolor{darkgray}{RGB}{55,55,55}
\definecolor{accentpurple}{RGB}{103,58,183}

\pagestyle{fancy}
\fancyhf{}
\rhead{EdgeLeak AI -- Informe Ejecutivo}
\lhead{Estado del Arte y Glosario Técnico}
\cfoot{\thepage}

\hypersetup{
  colorlinks=true,
  linkcolor=primaryblue,
  urlcolor=primaryblue,
  citecolor=primaryblue,
  pdftitle={Estado del Arte y Glosario Técnico - EdgeLeak AI},
  pdfauthor={Proyecto EdgeLeak AI},
}

\titleformat{\section}{\Large\bfseries\color{primaryblue}}{}{0em}{\thesection.\enspace}
\titleformat{\subsection}{\large\bfseries\color{darkgray}}{}{0em}{\thesubsection.\enspace}
\titleformat{\subsubsection}{\normalsize\bfseries}{}{0em}{\thesubsubsection.\enspace}

% ====================================================================
\title{
  \Huge\textbf{Informe Ejecutivo -- Sección de Estado del Arte}\\[0.5em]
  \textcolor{primaryblue}{\rule{\linewidth}{1.5pt}}\\[0.3em]
  \Large\textbf{EdgeLeak AI}\\[0.2em]
  \large Análisis de Brechas, Justificación Técnica y Glosario\\[0.5em]
  \normalsize Basado en el análisis de 13 referentes del estado del arte
}
\author{Proyecto EdgeLeak AI}
\date{\today}
% ====================================================================

\begin{document}

\maketitle
\thispagestyle{fancy}

\begin{mdframed}[linecolor=primaryblue,linewidth=1.5pt,backgroundcolor=blue!3]
\textbf{Alcance de esta sección:} Este documento constituye la sección de Estado del Arte
del informe ejecutivo de EdgeLeak~AI. Identifica las tres brechas principales detectadas
en la revisión sistemática del mapa de tendencias, presenta la justificación técnica de
por qué ninguno de los 13~referentes analizados resuelve dichas brechas de forma integral,
y concluye con un glosario completo de los términos técnicos empleados a lo largo del proyecto.
\end{mdframed}

\vspace{1em}
\tableofcontents
\newpage

% ====================================================================
\section{Estado del Arte: Análisis de 13 Referentes}
% ====================================================================

La revisión sistemática de la literatura y de productos comerciales arroja un universo
de 13~referentes organizados en cuatro categorías: trabajos académicos/prototipos (5),
papers científicos de alto impacto (5), documentación técnica oficial (1) y repositorios
\textit{open source} (1). La siguiente tabla consolida sus características principales y
las brechas que cada uno deja abiertas.

\begin{longtable}{>{\small}p{0.4cm} >{\small}p{3.8cm} >{\small}p{2.4cm} >{\small}p{2.4cm} >{\small}p{4.2cm}}
\caption{Matriz comparativa de los 13 referentes del estado del arte} \label{tab:matriz-comparativa} \\
\toprule
\textbf{\#} & \textbf{Referente} & \textbf{Hardware / Sensores} & \textbf{Software / IA} & \textbf{Brecha detectada} \\
\midrule
\endfirsthead
\multicolumn{5}{c}{\small\textit{(continúa en la siguiente página)}} \\
\toprule
\textbf{\#} & \textbf{Referente} & \textbf{Hardware / Sensores} & \textbf{Software / IA} & \textbf{Brecha detectada} \\
\midrule
\endhead
\bottomrule
\endlastfoot
1 & Anónimo (s.f.) -- Prototipo detección doméstica & Arduino Uno, sensor de humedad/nivel & Lógica condicional (If/Else), buzzer & Solo alerta ante inundación visible; no detecta microfugas ocultas. \\[0.4em]
2 & De Lima Rosado (2025) -- Seguridad IoT & N/A (arquitectura de red) & Protocolos de seguridad, cifrado WiFi & Enfoque exclusivo en ciberseguridad; sin innovación en detección física o Edge Computing. \\[0.4em]
3 & Islam \textit{et al.} (2023) -- IoT + ML, IEEE & Microcontroladores WiFi, medidores de flujo & ML clásico (Random Forest/SVM), Nube & Dependencia de la nube para ejecutar IA; latencia e inoperatividad sin internet. \\[0.4em]
4 & Kang \textit{et al.} (2023) -- Vibración, MDPI & Sensores piezoeléctricos de alta frecuencia & Análisis espectral, ML & Validado únicamente en entornos industriales; falla ante ruido doméstico cotidiano. \\[0.4em]
5 & Edge Impulse (2023-2024) -- TinyML/ESP32 & ESP32, micrófonos, acelerómetros & TinyML, Edge Impulse Studio, C++ & Herramienta genérica (vibración de motores); no calibrada para dinámica de fluidos en el hogar. \\[0.4em]
6 & Escobar \& López (2016) -- Redes de distribución & Arduino, sensor de flujo YF-S201 & Lógica en C++, pantalla LCD & El sensor de flujo tipo turbina no responde a micro-caudales (\textless{}0.1 L/min). \\[0.4em]
7 & Rubiños \textit{et al.} (2023) -- Smart Cities & Geófonos, medidores electromagnéticos macro & IA (redes neuronales, SVM) en servidores & Escala macro (acueductos); requiere alto cómputo en servidor; no aplicable al hogar. \\[0.4em]
8 & Farah \& Shahrour (2024) -- Review MDPI & Sensores acústicos, medidores IoT & IA, ML clásico & Foco en redes macro de empresas públicas; vacío total en Smart Home residencial. \\[0.4em]
9 & Corado (2025) -- Smart Cities, UOC & IoT, sensórica en tuberías & ML, modelos de optimización & Solución para infraestructura pública con servidores externos; sin Edge Computing. \\[0.4em]
10 & Islam \textit{et al.} (2022) -- Review IEEE & WSN, sensores de presión y acústicos & IoT, análisis de datos en la nube & Tuberías subterráneas y grandes caudales; sin filtrado de ruido doméstico. \\[0.4em]
11 & Sagat (s.f.) -- ESP8266 Leak Sensor (GitHub) & ESP8266, sensor genérico de lluvia & Arduino IDE, API PushBullet & Sensor pasivo (reactivo); sin IA ni sensor preventivo acústico. \\[0.4em]
12 & Dixit \textit{et al.} (2017) -- Calidad/Nivel IRJET & Sensores de presión, nivel y calidad & Notificaciones SMS & Depende de caída abrupta de presión; insensible a micro-goteos por goteo lento. \\[0.4em]
13 & Bruno \textit{et al.} (2021) -- Monitoreo presión MDPI & Arduino, transductores de presión, módulo SD & Datalogging pasivo & Sistema de registro diferido (SD); no alerta en tiempo real ni usa IA local. \\
\end{longtable}

\newpage

% ====================================================================
\section{Las Tres Brechas Principales del Estado del Arte}
% ====================================================================

Del análisis cruzado del mapa de tendencias y la revisión de los 13 referentes emergen
tres brechas estructurales que ningún referente resuelve de manera integral. Estas brechas
constituyen la \textbf{oportunidad de mercado y la justificación técnica} del proyecto
EdgeLeak~AI.

% -----------------------------------------------------------
\subsection{Brecha 1: Prevención y Escala --- Reactivo/Macro vs.\ Predictivo/Micro}
% -----------------------------------------------------------

\begin{mdframed}[linecolor=criticalred,linewidth=1pt,backgroundcolor=red!3]
\textbf{Hallazgo:} El mercado doméstico solo ofrece sensores reactivos de inundación
(detectan agua ya acumulada), mientras que las soluciones con Inteligencia Artificial
predictiva existen únicamente a escala de \emph{Smart City} para macro-acueductos.
Existe un vacío absoluto en la franja \emph{Smart Home} para micro-goteos domésticos
progresivos (rango de 0{,}1 a 1{,}0 L/min).
\end{mdframed}

\vspace{0.6em}

\subsubsection{Evidencia en los referentes}

Los sensores comerciales más vendidos (Xiaomi/Tuya Water Leak Sensor, referente~11
de Sagat) operan bajo un principio \textit{100\% reactivo}: solo generan una alarma
cuando el agua toca físicamente sus electrodos. Esto implica que la fuga ya es visible
y el daño al mobiliario ya ha comenzado. Por otra parte, las soluciones con IA genuina
---Corado (2025), Rubiños \textit{et al.} (2023), Farah \& Shahrour (2024), Islam
\textit{et al.} (2022)--- están diseñadas para monitorear tuberías de distribución
municipal (acueductos de ciudad), requiriendo sensores industriales costosos y capacidad
de cómputo en servidores externos.

\subsubsection{Referentes que NO resuelven esta brecha}

\begin{itemize}[nosep]
  \item Referentes 1 y 11 (Anónimo; Sagat): Reactivos por diseño, no predictivos.
  \item Referentes 7, 8 y 9 (Rubiños; Farah; Corado): Escala macro, incompatibles con domótica.
  \item Referentes 3, 10 (Islam 2022, 2023): Infraestructura subterránea, no residencial.
\end{itemize}

\subsubsection{Cómo EdgeLeak AI resuelve esta brecha}

EdgeLeak~AI traduce la inteligencia de detección de redes macro en un nodo compacto de
\textit{Smart Home}. Mediante un modelo \textbf{TinyML} entrenado con firmas acústicas
de goteo real y desplegado localmente en un \textbf{ESP32}, el sistema analiza
continuamente la vibración de la tubería bajo el lavaplatos. Esto permite anticipar la
microfuga en su fase de inicio ---cuando el desgaste del empaque genera turbulencia
audible pero el agua aún no ha llegado a acumularse--- convirtiendo un paradigma
\textit{reactivo} en uno genuinamente \textit{predictivo} para el entorno doméstico.

% -----------------------------------------------------------
\subsection{Brecha 2: Falsos Positivos --- Ausencia de Sensor Fusion}
% -----------------------------------------------------------

\begin{mdframed}[linecolor=warningorange,linewidth=1pt,backgroundcolor=orange!3]
\textbf{Hallazgo:} Los sensores acústicos y piezoeléctricos, aunque confirman ser el
estándar tecnológico para detectar fugas invisibles, producen un alto índice de
falsos positivos en entornos domésticos. El ruido generado por el uso cotidiano del
fregadero (lavado de platos, paso de agua) es indistinguible para un sistema monomodal
que analiza únicamente la señal acústica.
\end{mdframed}

\vspace{0.6em}

\subsubsection{Evidencia en los referentes}

Kang \textit{et al.} (2023) demuestran que las firmas de vibración son el método de
mayor precisión para detectar fugas presurizadas, pero sus pruebas se realizan en
entornos industriales controlados sin interferencias acústicas. Islam \textit{et al.}
(2022) reconocen explícitamente en su \textit{review} que los sensores acústicos en
tuberías subterráneas captan interferencias del entorno. El prototipo de Sagat (referente
11) y el sistema de Anónimo (referente 1) no cuentan con ningún mecanismo de
discriminación contextual: cualquier humedad activa la alarma, independientemente de
si es uso normal o fuga real.

\subsubsection{Referentes que NO resuelven esta brecha}

\begin{itemize}[nosep]
  \item Referentes 4 y 5 (Kang; Edge Impulse): Sensor único (vibración/audio); sin contexto de flujo.
  \item Referentes 1, 6 y 11 (Anónimo; Escobar; Sagat): Lógica binaria sin cruce de variables.
  \item Referentes 10 y 12 (Islam 2022; Dixit): No contemplan el ruido doméstico en su diseño.
\end{itemize}

\subsubsection{Cómo EdgeLeak AI resuelve esta brecha}

EdgeLeak~AI implementa una estrategia de \textbf{Sensor Fusion} que cruza dos señales
físicas complementarias:

\begin{enumerate}[nosep]
  \item \textbf{Señal acústica/vibración} (micrófono MEMS INMP441 o disco piezoeléctrico
    adherido al tubo): captura la firma espectral del fluido.
  \item \textbf{Señal de flujo} (sensor de efecto Hall YF-S201): actúa como
    \emph{interruptor de contexto}.
\end{enumerate}

\noindent La lógica de inferencia del modelo TinyML opera sobre la combinación de ambas
señales:

\begin{center}
\begin{tabular}{ll l}
\toprule
\textbf{Flujo} & \textbf{Señal acústica} & \textbf{Decisión del sistema} \\
\midrule
Caudal $>$ 0 & Ruido presente & Uso normal (grifo abierto) $\rightarrow$ Sin alerta \\
Caudal $= 0$ & Firma de goteo & Microfuga silenciosa $\rightarrow$ \textbf{Alerta crítica} \\
Caudal $= 0$ & Sin señal & Reposo normal $\rightarrow$ Sin alerta \\
\bottomrule
\end{tabular}
\end{center}

\noindent Este mecanismo de contexto elimina estructuralmente la categoría de falsos
positivos por uso cotidiano, problema que ningún referente del estado del arte resuelve.

% -----------------------------------------------------------
\subsection{Brecha 3: Privacidad y Eficiencia Energética --- Cloud vs.\ Edge AI}
% -----------------------------------------------------------

\begin{mdframed}[linecolor=accentpurple,linewidth=1pt,backgroundcolor=violet!3]
\textbf{Hallazgo:} La gran mayoría de los sistemas IoT con IA del estado del arte envían
datos continuamente a servidores externos (Cloud Computing) para ejecutar sus modelos.
Esto genera: (a)~alta latencia en la respuesta ante una fuga; (b)~dependencia de
conectividad permanente a internet; (c)~vulnerabilidades de privacidad al transmitir
audio del interior de un hogar a servidores de terceros; y (d)~consumo excesivo de ancho
de banda.
\end{mdframed}

\vspace{0.6em}

\subsubsection{Evidencia en los referentes}

Islam \textit{et al.} (2023) y Corado (2025) son representativos del paradigma
\textit{Cloud-first}: sus modelos de ML son ejecutados en servidores externos, lo que
implica que si el router falla, el sistema de detección queda ciego. De Lima Rosado
(2025) aborda específicamente las vulnerabilidades de seguridad que esta arquitectura
introduce, evidenciando que la transmisión continua de datos sensoriales desde el hogar
representa un vector de ataque. Bruno \textit{et al.} (2021) proponen un
\textit{datalogger} en tarjeta SD ---un sistema completamente desconectado que analiza
los datos de forma diferida---, evidenciando el extremo opuesto del problema: sin
inteligencia en tiempo real.

\subsubsection{Referentes que NO resuelven esta brecha}

\begin{itemize}[nosep]
  \item Referentes 3 y 9 (Islam 2023; Corado): IA exclusivamente en la nube.
  \item Referentes 7 y 8 (Rubiños; Farah): Alta capacidad de cómputo en servidor requerida.
  \item Referente 2 (De Lima Rosado): Identifica el problema de privacidad pero no lo soluciona.
  \item Referente 13 (Bruno): Procesa en diferido (SD card); sin alerta en tiempo real.
\end{itemize}

\subsubsection{Cómo EdgeLeak AI resuelve esta brecha}

EdgeLeak~AI adopta un paradigma de \textbf{Edge AI} (inteligencia en el borde): el
modelo TinyML reside y se ejecuta íntegramente dentro del microcontrolador ESP32.
La arquitectura de datos resultante es la siguiente:

\begin{enumerate}[nosep]
  \item \textbf{Captura}: Los sensores generan señales crudas localmente.
  \item \textbf{Inferencia local}: El ESP32 clasifica la señal en tiempo real
    (\textasciitilde{}ms de latencia) sin enviar audio ni datos crudos al exterior.
  \item \textbf{Transmisión mínima}: Solo el \textit{estado de alerta} (un mensaje JSON
    liviano) viaja por WiFi hacia la API de la aplicación Flutter.
  \item \textbf{Visualización}: La app móvil presenta el estado y el historial al usuario.
\end{enumerate}

Esta arquitectura garantiza: \textbf{baja latencia} (la inferencia no depende de la red),
\textbf{privacidad total} (el audio del hogar nunca abandona el dispositivo), y
\textbf{resiliencia} (funciona aunque el router pierda conectividad a internet).

\newpage

% ====================================================================
\section{Justificación Técnica Integral de EdgeLeak AI}
% ====================================================================

La siguiente tabla sintetiza, para cada uno de los 13~referentes, la razón específica
por la que su arquitectura no da respuesta al problema de las microfugas domésticas en
el punto de instalación bajo el lavaplatos, y el mecanismo concreto de EdgeLeak~AI que
cierra esa brecha.

\begin{longtable}{>{\small}p{0.4cm} >{\small}p{3.4cm} >{\small}p{5.0cm} >{\small}p{5.2cm}}
\caption{Justificación técnica: brecha por referente y respuesta de EdgeLeak AI} \label{tab:justificacion} \\
\toprule
\textbf{\#} & \textbf{Referente} & \textbf{Por qué NO resuelve el problema} & \textbf{Respuesta de EdgeLeak AI} \\
\midrule
\endfirsthead
\multicolumn{4}{c}{\small\textit{(continúa en la siguiente página)}} \\
\toprule
\textbf{\#} & \textbf{Referente} & \textbf{Por qué NO resuelve el problema} & \textbf{Respuesta de EdgeLeak AI} \\
\midrule
\endhead
\bottomrule
\endlastfoot
1 & Anónimo (s.f.) & Sistema binario y reactivo: solo alerta cuando el agua ya está estancada. No detecta microfugas ocultas. & TinyML analiza la vibración del tubo antes de que el agua sea visible; paradigma predictivo, no reactivo. \\[0.4em]
2 & De Lima Rosado (2025) & Enfoque estrictamente de ciberseguridad; sin innovación en detección física o procesamiento en el borde. & Edge AI garantiza privacidad al mantener el audio del hogar dentro del dispositivo; sin transmisión de datos crudos. \\[0.4em]
3 & Islam \textit{et al.} (2023) & Dependencia de la nube para ejecutar IA; latencia alta e inoperatividad sin internet continuo. & Inferencia TinyML local en ESP32: latencia de milisegundos y funcionalidad offline. \\[0.4em]
4 & Kang \textit{et al.} (2023) & Validado en entornos industriales controlados; genera falsos positivos ante el ruido doméstico. & Sensor Fusion (flujo + vibración) proporciona contexto que discrimina uso normal de fuga real. \\[0.4em]
5 & Edge Impulse (2023-2024) & Herramienta genérica para vibración de motores; no calibrada para la dinámica física del goteo en tuberías domésticas. & Modelo TinyML entrenado específicamente con firmas acústicas de microfugas en tubería de lavaplatos. \\[0.4em]
6 & Escobar \& López (2016) & Sensor YF-S201 solo detecta caudales $>$1~L/min; insensible a micro-goteos de 0.1--0.5~L/min. & El sensor de flujo se complementa con el sensor acústico; el micro-goteo es detectado por la firma vibratoria, no por la turbina. \\[0.4em]
7 & Rubiños \textit{et al.} (2023) & Escala macro (acueductos urbanos); requiere servidores de alto cómputo incompatibles con domótica. & Miniaturización de la IA mediante TinyML: mismo principio de redes neuronales en un ESP32 de \$5~USD. \\[0.4em]
8 & Farah \& Shahrour (2024) & Foco en redes macro de empresas de servicios públicos; vacío total en el ecosistema Smart Home. & Primer nodo doméstico que aplica análisis acústico con IA al punto de falla más común en el hogar (el latiguillo). \\[0.4em]
9 & Corado (2025) & Solución para infraestructura pública con servidores externos; sin Edge Computing ni aplicación residencial. & Arquitectura Edge-first: toda la inteligencia vive en el dispositivo; la nube solo recibe el estado de alerta. \\[0.4em]
10 & Islam \textit{et al.} (2022) & Tuberías subterráneas y grandes caudales; sin mecanismo para filtrar interferencias de una cocina residencial. & Sensor Fusion como filtro de contexto: el sensor de flujo invalida la alarma acústica cuando el grifo está abierto. \\[0.4em]
11 & Sagat (s.f.) & Sensor pasivo (requiere que el agua moje los pines); sin IA, sin análisis de señal preventivo. & Micrófono MEMS/disco piezoeléctrico activo: escucha la tubería en tiempo real sin necesidad de contacto con el agua. \\[0.4em]
12 & Dixit \textit{et al.} (2017) & Depende de caída abrupta de presión, método ineficaz para micro-goteos donde la presión casi no varía. & La firma acústica de la turbulencia del goteo es independiente de la presión de red; detecta el fenómeno físico real. \\[0.4em]
13 & Bruno \textit{et al.} (2021) & Sistema de registro diferido (SD card); sin alerta en tiempo real ni inteligencia local. & El ESP32 reemplaza la tarjeta SD por un modelo TinyML que procesa la señal instantáneamente y notifica al móvil en tiempo real. \\
\end{longtable}

\begin{mdframed}[linecolor=normalgreen,linewidth=1.5pt,backgroundcolor=green!3]
\textbf{Síntesis:} EdgeLeak~AI es el único sistema, entre todos los referentes analizados,
que combina simultáneamente: (1)~\textbf{Prevención activa} mediante análisis de firma
acústica, (2)~\textbf{Sensor Fusion} para eliminar falsos positivos en el contexto
doméstico, y (3)~\textbf{Edge AI} para garantizar privacidad, baja latencia y resiliencia
ante cortes de internet. Esta combinación lo posiciona como una solución inédita en el
nicho de las microfugas en instalaciones domésticas bajo el fregadero.
\end{mdframed}

\newpage

% ====================================================================
\section{Glosario de Términos Técnicos}
% ====================================================================

El presente glosario reúne los términos especializados empleados a lo largo de la
documentación técnica y el informe ejecutivo de EdgeLeak~AI, organizados alfabéticamente
para facilitar su consulta.

\begin{longtable}{>{\bfseries}p{3.8cm} p{10.5cm}}
\caption{Glosario de términos técnicos -- EdgeLeak AI} \label{tab:glosario} \\
\toprule
\textbf{Término} & \textbf{Definición} \\
\midrule
\endfirsthead
\multicolumn{2}{c}{\small\textit{(continúa en la siguiente página)}} \\
\toprule
\textbf{Término} & \textbf{Definición} \\
\midrule
\endhead
\bottomrule
\endlastfoot

ADC (Analog-to-Digital Converter) & Conversor analógico-digital. Circuito que transforma
una señal eléctrica continua (analógica) en un valor numérico discreto procesable por un
microcontrolador. En el proyecto se evalúa el módulo externo ADS1115 (16 bits) para
superar la imprecisión del ADC interno del ESP32 al medir vibraciones milimétricas.
\\[0.5em]

Anomaly Detection & Detección de anomalías. Técnica de Machine Learning no supervisado
que identifica patrones que se desvían significativamente de la línea base establecida.
Edge Impulse la implementa como bloque de aprendizaje en la pipeline de TinyML para
detectar la firma acústica del goteo como una anomalía frente al reposo del tubo.
\\[0.5em]

API (Application Programming Interface) & Interfaz de Programación de Aplicaciones.
Conjunto de protocolos y definiciones que permiten la comunicación entre sistemas de
software. En EdgeLeak~AI, la API de Groq expone el modelo LLM para el análisis de
patrones desde la aplicación Flutter.
\\[0.5em]

Arduino & Plataforma de hardware y software de código abierto basada en microcontroladores
AVR/ARM. Referenciada en múltiples trabajos del estado del arte (Anónimo, s.f.; Escobar
\& López, 2016; Bruno \textit{et al.}, 2021) como base de prototipos IoT de bajo costo.
\\[0.5em]

Async/Await & Modelo de programación asíncrona del lenguaje Dart (y otros lenguajes
modernos) que permite ejecutar operaciones no bloqueantes (llamadas de red, acceso a
BD) sin congelar la interfaz de usuario. Fundamental en la arquitectura de los
controllers de EdgeLeak~AI.
\\[0.5em]

ChangeNotifier & Clase del framework Flutter que implementa el patrón Observer. Los
controladores de EdgeLeak~AI la extienden para notificar reactivamente a los widgets
cuando el estado cambia (nueva lectura del sensor, nueva alerta, cambio de sesión).
\\[0.5em]

Cloud Computing & Paradigma de computación en el que el procesamiento y almacenamiento
de datos ocurre en servidores remotos accesibles por internet. Contraposición al modelo
Edge Computing. La mayoría de los referentes del estado del arte (Islam, 2023; Corado,
2025) dependen de este paradigma para ejecutar sus modelos de IA.
\\[0.5em]

CRUD & Acrónimo de las cuatro operaciones básicas sobre datos persistentes:
\textbf{C}reate (Crear), \textbf{R}ead (Leer), \textbf{U}pdate (Actualizar),
\textbf{D}elete (Eliminar). El módulo de administración de EdgeLeak~AI implementa
CRUD completo sobre la entidad \texttt{UsuarioModel}.
\\[0.5em]

Dart & Lenguaje de programación de tipado estático y fuerte desarrollado por Google.
Es el lenguaje base del framework Flutter. Soporta null-safety, programación orientada
a objetos y paradigmas asíncronos mediante async/await.
\\[0.5em]

Datalogging & Registro diferido de datos sensoriales en un medio de almacenamiento local
(típicamente tarjeta SD o memoria Flash) para su análisis posterior. Representado en el
estado del arte por Bruno \textit{et al.} (2021). EdgeLeak~AI supera este paradigma
mediante inferencia TinyML en tiempo real.
\\[0.5em]

Edge AI & Inteligencia Artificial desplegada directamente en el dispositivo de campo
(el \textit{``borde''} de la red), sin dependencia de servidores remotos para la
inferencia. Permite baja latencia, privacidad de datos y funcionamiento offline.
Principio arquitectónico central de EdgeLeak~AI.
\\[0.5em]

Edge Computing & Paradigma de computación distribuida en el que el procesamiento de datos
se realiza en el propio dispositivo o en nodos cercanos a la fuente de datos, en lugar
de en la nube central. Contrastado con Cloud Computing en los referentes 3, 7 y 9 del
estado del arte.
\\[0.5em]

Edge Impulse & Plataforma de desarrollo de TinyML (desarrollada por Edge Impulse Inc.)
que provee un entorno integrado para recolección de datos, entrenamiento de modelos de
IA y generación de firmware optimizado para microcontroladores como el ESP32. Utilizada
como referente técnico (referente~5) y como herramienta metodológica del proyecto.
\\[0.5em]

ESP32 & Microcontrolador de doble núcleo (Xtensa LX6/LX7) con WiFi y Bluetooth
integrados, fabricado por Espressif Systems. Es el \textit{``cerebro''} del hardware
de EdgeLeak~AI. Su variante DEVKIT V1 (WROOM-32) cuesta aproximadamente \$4--6~USD y
es 100\% compatible con Edge Impulse para el despliegue de modelos TinyML.
\\[0.5em]

ESP8266 & Predecesor del ESP32; microcontrolador con WiFi integrado de un solo núcleo.
Empleado en el referente~11 (Sagat, s.f.) como base del sistema de notificaciones Push
de bajo costo. El ESP32 lo supera en capacidad de cómputo y soporte para IA.
\\[0.5em]

Efecto Hall & Principio físico por el cual un conductor que transporta corriente eléctrica
en presencia de un campo magnético perpendicular genera un voltaje transversal medible.
Los sensores de flujo YF-S201/YF-B1 utilizan este efecto para contar las rotaciones de
una turbina y calcular el caudal de agua en litros por minuto.
\\[0.5em]

Falso Positivo & En clasificación de Machine Learning, evento en el que el sistema predice
la clase positiva (ej.\ ``fuga detectada'') cuando la realidad es negativa (ej.\ uso
normal del grifo). La Brecha~2 del estado del arte identifica la ausencia de Sensor
Fusion como causa principal de falsos positivos en sistemas de detección acústica
doméstica.
\\[0.5em]

Firma Acústica & Patrón característico de una señal de audio o vibración que permite
identificar la fuente que la generó. Cada tipo de fuga (goteo lento, presurizado, sifón)
genera una firma acústica única en la tubería. El modelo TinyML de EdgeLeak~AI aprende
a reconocer estas firmas durante la fase de entrenamiento.
\\[0.5em]

Flutter & Framework de desarrollo de aplicaciones móviles, web y de escritorio de código
abierto, desarrollado por Google. Utiliza el lenguaje Dart y proporciona un motor de
renderizado propio (Skia/Impeller). Es el framework de la capa de presentación (Vista)
de EdgeLeak~AI.
\\[0.5em]

Golpe de Ariete & Fenómeno hidráulico que ocurre cuando el flujo de un fluido en una
tubería se detiene o cambia bruscamente de dirección (al cerrar rápidamente una válvula
o grifo), generando una onda de presión que puede dañar acoples y empaques con el tiempo.
Es uno de los mecanismos de desgaste que origina las microfugas que detecta EdgeLeak~AI.
\\[0.5em]

Groq API & Interfaz de programación de aplicaciones del proveedor de infraestructura de
IA Groq, que permite invocar modelos de lenguaje de gran escala (LLM) de alta velocidad
a través de llamadas REST. EdgeLeak~AI utiliza el modelo \texttt{Llama-3.3-70B-Versatile}
a través de esta API para el análisis avanzado de patrones en la fase sin hardware físico.
\\[0.5em]

I2S (Inter-IC Sound) & Protocolo de comunicación serie sincrónico estándar para
transmisión de datos de audio digital entre circuitos integrados. El micrófono MEMS
INMP441 utiliza I2S para comunicarse con el ESP32, entregando señales de audio digitales
de alta calidad sin el ruido inherente a los conversores ADC analógicos.
\\[0.5em]

INMP441 & Módulo de micrófono MEMS (Micro-Electro-Mechanical Systems) omnidireccional
con interfaz digital I2S. Opción de sensor acústico recomendada para el hardware de
EdgeLeak~AI por su alta fidelidad de señal y compatibilidad directa con el ESP32 para
captura de datos de entrenamiento TinyML.
\\[0.5em]

Inferencia (TinyML) & Proceso mediante el cual un modelo de Machine Learning ya entrenado
analiza datos nuevos y produce una predicción o clasificación. En el contexto de TinyML,
la inferencia ocurre directamente en el microcontrolador (ESP32) sin necesidad de enviar
los datos a la nube.
\\[0.5em]

IoT (Internet of Things) & Internet de las Cosas. Paradigma tecnológico que conecta
dispositivos físicos del mundo real (sensores, actuadores, electrodomésticos) a internet
para recopilar datos y automatizar acciones. EdgeLeak~AI es un nodo IoT de Smart Home.
\\[0.5em]

IP65 / IP67 & Clasificaciones del estándar internacional IEC 60529 que describen el
grado de protección de una carcasa contra partículas sólidas y líquidos. IP65 garantiza
protección total contra polvo y resistencia a chorros de agua; IP67 añade resistencia a
inmersión temporal. Requisito para la caja electrónica del hardware de EdgeLeak~AI dado
el entorno de humedad bajo el lavaplatos.
\\[0.5em]

JSON (JavaScript Object Notation) & Formato ligero de intercambio de datos basado en
texto, legible por humanos y de fácil parseo por máquinas. Utilizado en EdgeLeak~AI para
la serialización/deserialización de modelos de datos y para los mensajes de estado
transmitidos por WiFi entre el ESP32 y la aplicación Flutter.
\\[0.5em]

Latiguillo & Manguera flexible de conexión hidráulica, generalmente de acero inoxidable
trenzado con terminales roscados, usada para conectar la tubería de agua fría a la
llave del fregadero. Es el componente de mayor frecuencia de fallo en la instalación
bajo el lavaplatos y el punto de instalación primario del sensor de EdgeLeak~AI.
\\[0.5em]

LLM (Large Language Model) & Modelo de Lenguaje de Gran Escala. Tipo de modelo de
Inteligencia Artificial entrenado sobre enormes corpus de texto que puede comprender y
generar lenguaje natural, razonar y analizar datos. EdgeLeak~AI utiliza
\texttt{Llama-3.3-70B-Versatile} (vía API Groq) como motor de análisis de patrones
en ausencia de hardware físico.
\\[0.5em]

Machine Learning (ML) & Subcampo de la Inteligencia Artificial que desarrolla algoritmos
capaces de aprender patrones a partir de datos y mejorar su rendimiento con la
experiencia, sin ser explícitamente programados para cada caso. Incluye técnicas
supervisadas (SVM, Random Forest) y no supervisadas (detección de anomalías).
\\[0.5em]

MEMS (Micro-Electro-Mechanical Systems) & Sistemas microelectromecánicos. Tecnología de
fabricación que integra componentes mecánicos, sensores y electrónica en chips de silicio
de escala micrométrica. Los micrófonos MEMS (como el INMP441) ofrecen alta calidad de
señal en formatos miniaturizados, ideales para dispositivos IoT.
\\[0.5em]

Microfuga & Fuga de agua de caudal muy bajo (rango de 0.1 a 1.0~L/min) generada por el
desgaste progresivo de empaques, roscas o acoples. Por su bajo caudal, no es registrada
por los medidores de agua convencionales (zona muerta mecánica) y no genera charcos
visibles hasta semanas después de su inicio. Es el tipo de fallo objetivo de EdgeLeak~AI.
\\[0.5em]

MVC (Model-View-Controller) & Patrón arquitectónico de software que separa la aplicación
en tres capas: \textbf{Modelo} (datos y lógica de negocio), \textbf{Vista} (presentación)
y \textbf{Controlador} (mediación entre modelo y vista). EdgeLeak~AI implementa MVC
adaptado a las convenciones reactivas de Flutter mediante \texttt{ChangeNotifier}.
\\[0.5em]

Null-Safety & Característica del lenguaje Dart que garantiza en tiempo de compilación
que ninguna variable puede contener un valor nulo salvo que sea declarada explícitamente
como nullable (\texttt{?}). Previene una categoría entera de errores de ejecución
(NullPointerException).
\\[0.5em]

Piezoeléctrico & Material o componente que genera una tensión eléctrica cuando se le
aplica una deformación mecánica (presión, vibración), y viceversa. Los discos
piezoeléctricos son la opción de sensor de vibración de menor costo (\textless{}\$1~USD)
para capturar la firma acústica del goteo en la tubería del lavaplatos.
\\[0.5em]

Random Forest & Algoritmo de Machine Learning supervisado que construye un conjunto
(\textit{ensemble}) de árboles de decisión y combina sus predicciones para mejorar la
precisión y evitar el sobreajuste. Citado por Islam \textit{et al.} (2023) como uno de
los algoritmos de ML clásico aplicados a la detección de fugas.
\\[0.5em]

Red Neuronal & Modelo matemático de Machine Learning compuesto por capas de nodos
(neuronas artificiales) interconectadas con pesos ajustables. Las redes neuronales
convolucionales (CNN) y densas son los tipos utilizados por Edge Impulse para clasificar
señales de audio/vibración en modelos TinyML.
\\[0.5em]

Sensor Fusion & Fusión de sensores. Técnica que combina datos de múltiples sensores
heterogéneos para producir información más precisa y contextual que la que puede proveer
cada sensor individualmente. EdgeLeak~AI aplica Sensor Fusion cruzando la señal acústica
con la señal de flujo para eliminar falsos positivos.
\\[0.5em]

Sifón & Componente del desagüe del fregadero con forma de ``U'' o ``S'' que retiene agua
para impedir la entrada de gases malolientes del alcantarillado. Es un segundo punto de
posible fuga en la instalación bajo el lavaplatos.
\\[0.5em]

Smart City & Ciudad Inteligente. Concepto urbano que integra tecnologías IoT y análisis
de datos a gran escala para optimizar servicios públicos (agua, energía, transporte).
La mayoría de las soluciones de detección de fugas con IA del estado del arte pertenecen
a este ecosistema macro, inaccesible para el usuario residencial.
\\[0.5em]

Smart Home & Hogar Inteligente. Ecosistema doméstico en el que dispositivos y sensores
están interconectados para automatizar, monitorear y mejorar la eficiencia y seguridad
del hogar. EdgeLeak~AI es un nodo Smart Home de bajo costo para el nicho de las
microfugas domésticas.
\\[0.5em]

SMTP (Simple Mail Transfer Protocol) & Protocolo estándar de internet para el envío de
correo electrónico. EdgeLeak~AI utiliza Google Workspace SMTP (con SSL/TLS) a través del
servicio \texttt{EmailService} para enviar notificaciones y correos de recuperación de
contraseña.
\\[0.5em]

SQLite & Motor de base de datos relacional integrado, ligero y sin servidor. EdgeLeak~AI
lo utiliza para la persistencia local de datos en dispositivos Android/iOS mediante el
paquete \texttt{sqflite}, almacenando usuarios, alertas y lecturas de sensores en el
archivo \texttt{edgeleak\_v3.db}.
\\[0.5em]

SSL/TLS (Secure Sockets Layer / Transport Layer Security) & Protocolos criptográficos
que garantizan la confidencialidad, integridad y autenticidad de las comunicaciones en
redes de computadores. Utilizados en EdgeLeak~AI para asegurar la comunicación SMTP de
correo y las llamadas a la API de Groq.
\\[0.5em]

SVM (Support Vector Machine) & Máquina de Vectores de Soporte. Algoritmo de Machine
Learning supervisado que encuentra el hiperplano de máximo margen que separa las clases
en el espacio de características. Citado por Rubiños \textit{et al.} (2023) e Islam
\textit{et al.} (2023) como técnica clásica aplicada a la clasificación de patrones de
fugas.
\\[0.5em]

TDD (Test-Driven Development) & Desarrollo Guiado por Pruebas. Metodología de desarrollo
de software en la que se escriben primero las pruebas unitarias (que fallan), luego el
código mínimo necesario para que pasen y, finalmente, se refactoriza. EdgeLeak~AI adopta
TDD mediante el framework \texttt{flutter\_test}.
\\[0.5em]

TinyML & Subconjunto de Machine Learning orientado a ejecutar modelos de inferencia en
dispositivos de baja potencia y recursos limitados (microcontroladores, sensores
embebidos). Es la tecnología núcleo del hardware de EdgeLeak~AI, permitiendo ejecutar
la clasificación de firmas acústicas directamente en el ESP32.
\\[0.5em]

Turbulencia (fluidos) & Régimen de flujo de fluidos caracterizado por variaciones
caóticas de velocidad y presión (en contraposición al flujo laminar). Una microfuga
en un acople genera turbulencia localizada que produce ondas de vibración de alta
frecuencia detectables por el sensor acústico de EdgeLeak~AI antes de que el agua
sea visible.
\\[0.5em]

WSN (Wireless Sensor Network) & Red Inalámbrica de Sensores. Conjunto de nodos
distribuidos espacialmente que monitorean condiciones físicas y comunican los datos
de forma inalámbrica. Citado por Islam \textit{et al.} (2022) como topología de red
para monitoreo de tuberías subterráneas.
\\[0.5em]

WiFi (IEEE 802.11) & Estándar de comunicación inalámbrica de área local de alta
velocidad. Es el protocolo de comunicación utilizado por el ESP32 de EdgeLeak~AI para
transmitir el estado de alerta desde el nodo hardware hacia la aplicación Flutter a
través de la red doméstica.
\\[0.5em]

YF-S201 / YF-B1 & Sensores de flujo de agua de efecto Hall de bajo costo (plástico y
latón, respectivamente) con roscas de 1/2~pulgada. Miden el caudal en L/min mediante el
conteo de pulsos generados por una turbina interna. En EdgeLeak~AI actúan como
\textit{interruptores de contexto} para la estrategia de Sensor Fusion.
\\[0.5em]

Zona Muerta Mecánica & Rango de caudal por debajo del cual los medidores de agua
convencionales (tipo Woltmann o mecánicos) no registran flujo porque la turbina interna
no tiene energía suficiente para girar. Típicamente se sitúa en valores menores a
0.5~L/min, lo que hace que las microfugas sean invisibles para la facturación del
servicio público y para el usuario.
\\

\end{longtable}

% ====================================================================
\section*{Referencias Bibliográficas}
\addcontentsline{toc}{section}{Referencias Bibliográficas}
% ====================================================================

\begin{enumerate}[label={[\arabic*]}, leftmargin=*, nosep, itemsep=0.4em]

\item Anónimo (s.f.). \textit{Prototipo de un sistema o red de sensores para la detección
de fugas de agua de uso doméstico utilizando un microcontrolador arduino}. [Trabajo
académico / prototipo].

\item Bruno, F., De Marchis, M., Milici, B., Saccone, D., \& Traina, F. (2021).
A Pressure Monitoring System for Water Distribution Networks Based on Arduino
Microcontroller. \textit{Water, 13}(17), 2321.
\url{https://doi.org/10.3390/w13172321}

\item Corado Pérez, L. (2025). \textit{Detección y prevención de fugas en tuberías e
infraestructuras de agua} [Trabajo Fin de Máster, Universitat Oberta de Catalunya].
Repositorio Institucional UOC.

\item De Lima Rosado, E. O. (2025). \textit{Fortalecimiento de la seguridad en el
internet de las cosas (IoT): estrategias y propuestas para mitigar riesgos de
vulnerabilidad en sensores de agua con tecnología IoT} [Tesis de Especialización,
Universidad Nacional Abierta y a Distancia -- UNAD]. Repositorio Institucional UNAD.

\item Dixit, R., Chaudhari, H., Jadhav, S., \& Jagtap, K. (2017). Water Level and
Leakage Detection System with its Quality Analysis based on Sensor for Home Application.
\textit{International Research Journal of Engineering and Technology (IRJET), 4}(11).

\item Edge Impulse. (2023). \textit{Anomaly Detection Guide}. Documentación técnica
oficial. \url{https://docs.edgeimpulse.com/docs/edge-impulse-studio/learning-blocks/anomaly-detection}

\item Escobar Sánchez, J. A., \& López Hernández, F. (2016). \textit{Detección de fugas
en una red de distribución de agua potable} [Trabajo de Residencia Profesional,
Instituto Tecnológico de Tuxtla Gutiérrez]. Tecnológico Nacional de México.

\item Farah, E., \& Shahrour, I. (2024). Water Leak Detection: A Comprehensive Review of
Methods, Challenges, and Future Directions. \textit{Water, 16}(21), 2975.
\url{https://doi.org/10.3390/w16212975}

\item Islam, M. R., Azam, S., Shanmugam, B., \& Mathur, D. (2022). A Review on Current
Technologies and Future Direction of Water Leakage Detection in Water Distribution
Network. \textit{IEEE Access, 10}, 107177--107201.
\url{https://doi.org/10.1109/ACCESS.2022.3212769}

\item Islam, M. R., et al. (2023). An Intelligent IoT and ML-Based Water Leakage
Detection System. \textit{IEEE Access, 11}, 10305165.
\url{https://ieeexplore.ieee.org/document/10305165}

\item Kang, J., et al. (2023). Machine Learning Model for Leak Detection Using Water
Pipeline Vibration Sensor. \textit{Sensors, 23}(21), 8935.
\url{https://doi.org/10.3390/s23218935}

\item Rubiños Jiménez, S. L., Cuzcano Rivas, A. B., Tejada Cabanillas, A. A., Mendoza
Apaza, F., Leva Apaza, A., \& Apesteguia Infantes, J. A. (2023). \textit{Los sistemas de
abastecimiento de agua potable: prevención y detección temprana de fugas}. Centro de
Investigación y Desarrollo Ecuador (CIDE).
\url{https://doi.org/10.33936/978-9942-636-36-2}

\item Sagat, P. (s.f.). \textit{ESP8266 Nodemcu Leak Sensor with SMS functionality}
[Repositorio de código fuente]. GitHub.
\url{https://github.com/psagat/ESP8266-Leak-Sensor}

\end{enumerate}

\end{document}
